\chapter{Singular K3 surfaces and complex multiplication}\label{ch:singular}

Where the Picard number of a one-parameter family jumps from $19$ to $20$, the
transcendental lattice drops from rank $3$ to rank $2$: it becomes a
positive-definite even binary quadratic form, and by Shioda--Inose that form
determines the surface. In the $M_n$-polarized families these jumps happen at
the CM points of the modular curve $X_0(n)^+$, and the binary form is the
orthogonal complement of the algebraic class that appeared. This chapter is
where the lattice arithmetic of Chapter~\ref{ch:lattices} pays for itself.

\section{Shioda--Inose structures}

\begin{definition}[Morrison~\cite{Morrison1984}]
A K3 surface $X$ has a \emph{Shioda--Inose structure} if it carries a
symplectic involution $\iota$ such that the minimal resolution $Y$ of
$X/\iota$ is a Kummer surface $\Km(A)$, and the induced rational map
$X\dashrightarrow Y$ of degree $2$ induces a Hodge isometry
$T(X)\cong T(A)$ (with the form on $T(A)$ multiplied by $2$ on one side,
depending on conventions --- see the source).
\end{definition}

\begin{theorem}[\Lstatus, Morrison~\cite{Morrison1984}, Thm.~6.3]
A K3 surface $X$ admits a Shioda--Inose structure if and only if there is a
primitive embedding $T(X)\hookrightarrow\U^3$; equivalently, if
$E_8(-1)^2\hookrightarrow\NS(X)$ primitively. Every K3 surface with
$\rho\ge19$ has one.
\end{theorem}

For $\rho=19$ this is the statement that $T\cong\U\oplus\lat{2n}$ embeds in
$\U^3$ (Example~\ref{ex:Mn} exhibits it for $n=7$), and it is the
lattice-theoretic reason the $M_n$-polarized surfaces are ``sandwiched'' between
Kummer surfaces of products of isogenous elliptic curves. For $\rho=20$:

\begin{theorem}[\Lstatus, Shioda--Inose~\cite{ShiodaInose1977}]\label{thm:SI-construction}
Let $T$ be a positive-definite even binary form of discriminant $D>0$, with
associated imaginary quadratic order of discriminant $-D$, and let
$\tau_1,\tau_2\in\HH$ be such that $E_{\tau_1}\times E_{\tau_2}$ is the abelian
surface with $T(E_{\tau_1}\times E_{\tau_2})\cong T$ (this exists and is unique
up to isomorphism). Then the K3 surface with transcendental lattice $T$ is
the Shioda--Inose partner of $E_{\tau_1}\times E_{\tau_2}$: explicitly, the
Inose surface $X_{\alpha,\beta}$ of \S\ref{sec:inose} with $J_i=j(\tau_i)/1728$.
\end{theorem}

\begin{example}
$T=A_2$, $D=3$: $E_\omega\times E_\omega$, $J_1=J_2=0$, $(\alpha,\beta)=(0,\pm1)$.
$T=\lat2\oplus\lat2$, $D=4$: $E_i\times E_i$, $J_1=J_2=1$, $(\alpha,\beta)=(1,0)$
(with $\alpha^3=1$; the three cube roots give isomorphic surfaces). These are
the $X_3$ and $X_4$ of Example~\ref{ex:inose-points}.
\end{example}

\section{Rank jumps along a one-parameter family}

Let $X_\tau$, $\tau\in\HH$, be the $M_n$-polarized family with period
$\omega(\tau)=e-n\tau^2f+\tau w\in T_n\otimes\C$ (Chapter~\ref{ch:torelli}).
For $v\in T_n$ the \emph{wall} of $v$ is $\{\tau:\langle v,\omega(\tau)\rangle=0\}$;
by Lefschetz $(1,1)$, at such $\tau$ the class $v$ is algebraic, so
\[
  \NS(X_\tau)\supseteq M_n\oplus\Z v\quad\text{(finite index)},\qquad
  \rho(X_\tau)=20,\qquad
  T(X_\tau)=v^{\perp}\cap T_n \quad\text{(when $v$ is primitive and $\NS$ is exactly the saturation)}.
\]
Since $\langle v,\omega(\tau)\rangle=-nv_0\tau^2+v_1+2nv_2\tau$ is a quadratic
in $\tau$ with integer coefficients, its roots are quadratic imaginary
numbers when $v^2<0$: the wall of $v$ consists of \emph{CM points}. The
discriminant of the quadratic is $4n^2v_2^2+4nv_0v_1=2n\,v^2$, so
\[
  \tau=\frac{-2nv_2\pm\sqrt{2n\,v^2}}{-2nv_0}
  \quad\text{is a CM point of discriminant $2n\,v^2/\gcd^2$},
\]
and the transcendental lattice at that point is the binary form $v^{\perp}$,
whose discriminant is $\disc(v^{\perp})=-v^2\cdot 2n/\divv(v)^2$ by
Proposition~\ref{prop:index}. Everything in this paragraph is elementary
except the identification $T(X_\tau)=v^{\perp}$, which uses that $\NS(X_\tau)$
is exactly the saturation of $M_n\oplus\Z v$ (true for the generic point of
the wall by dimension count and Torelli) \Lstatus.

\begin{example}[the walls of $T_7$, revisited]\label{ex:walls-T7}
Take $n=7$ and the four classes of Chapter~\ref{ch:lattices}.
\begin{center}
\begin{tabular}{@{}lllll@{}}
\toprule
$v$ & $v^2$ & $\divv v$ & wall $\tau$ & $T(X_\tau)=v^{\perp}$, $\disc$\\
\midrule
$(1,-1,0)$ & $-2$ & $1$ & $7\tau^2=-1$: $\tau=i/\sqrt7$ & $\lat2\oplus\lat{14}$, $28$\\
$(2,-4,1)$ & $-2$ & $2$ & $14\tau^2-14\tau+4=0$: $\tau=\tfrac12+\tfrac{i}{2\sqrt7}$ & $\bigl(\begin{smallmatrix}2&1\\1&4\end{smallmatrix}\bigr)$, $7$\\
$(14,-14,-5)$ & $-42$ & $14$ & $\tau=-\tfrac5{14}+\tfrac{i\sqrt3}{14}$ & $A_2$, $3$\\
\bottomrule
\end{tabular}
\end{center}
The lattice column is kernel-checked \Kstatus\ (Examples~\ref{ex:fricke-root}--\ref{ex:disc3});
the wall column is the elementary quadratic above \Estatus. Note the
discriminants $28$, $7$, $3$: by Theorem~\ref{thm:shioda-inose} these three
points of the $M_7$-family are the singular K3 surfaces with those
transcendental lattices, and the third is $X_3$ itself. The CM points have
discriminants $-28$, $-7$, $-3$ as points of $\HH$ --- the discriminant of the
transcendental lattice and the discriminant of the CM point agree, as they must
(both are $2n\,v^2/\divv(v)^2$ up to sign).
\end{example}

\begin{example}[$T_{10}$]
For $n=10$: $(1,-1,0)$ gives $\lat2\oplus\lat{20}$ ($\disc40$) at
$\tau=i/\sqrt{10}$; $(10,-10,-3)$ gives $\lat2\oplus\lat2$ ($\disc4$) at
$\tau=-\tfrac3{10}+\tfrac i{10}$, i.e.\ $X_4$ \Kstatus\ for the lattice
(Example~\ref{ex:disc4}), \Estatus\ for the wall.
\end{example}

\begin{warning}[forced, not corroborating]\label{warn:forced}
When a program computes the CM point $\tau$ numerically, recognizes it as an
element of an imaginary quadratic field of discriminant $-D$, and separately
computes the binary form $v^{\perp}$ and finds its discriminant to be $D$, it
is tempting to report ``two independent computations agree''. They are not
independent: both are $2n\,v^2/\divv(v)^2$. The agreement is \emph{forced} by
the algebra above and carries no evidential weight about the surface. What
\emph{does} carry weight is whether the recognised $\tau$ maps, under the
family's parametrization $z(\tau)$, to a point of the $z$-line that one can
check independently --- for instance a singular point of the Picard--Fuchs
operator. Chapter~\ref{ch:frontier} reports such a check.
\end{warning}

\section{Which points are ample, and which are not}

Dolgachev's Theorem~\ref{thm:dolgachev-moduli}(iii) removes the $(-2)$-walls
from the moduli space of \emph{pseudo-ample} polarized surfaces: at the wall
of a root $\delta$, the class $\delta$ is effective and the polarization
class meets it trivially, so the polarization is no longer ample on the
surface; the ``surface at the wall'' is a member with an extra $(-2)$-curve,
or its contraction to a nodal surface. So among the three points of
Example~\ref{ex:walls-T7}, the first two (roots) are of that nature, while the
third (norm $-42$, not a root) is a genuine smooth $M_7$-polarized surface of
Picard number $20$ --- and it is the one with $T=A_2$, of smallest possible
discriminant. This is not a coincidence of the example but a general pattern:
the smallest-discriminant CM point of $X_0(n)^+$ is never on a $(-2)$-wall
when $n>1$, since a root's complement has discriminant $2n/\divv^2\ge 2n/4$.

\section{Exercises}

\begin{exercise}[$\star$]
Derive the discriminant $2n\,v^2$ of the wall quadratic and the formula
$\disc(v^{\perp})=-2n\,v^2/\divv(v)^2$, and conclude that the CM discriminant
and the form discriminant always agree.
\end{exercise}

\begin{exercise}
For $n=7$, enumerate all primitive $v\in T_7$ with $-v^2\le 44$ up to the
action of the swap and translation, compute $\disc(v^{\perp})$, and list the
distinct binary forms that occur. Which of them are on $(-2)$-walls?
\end{exercise}

\begin{exercise}
Show that the wall equations of $(2,-4,1)$ and $(-2,4,1)$ in $T_7$ have roots
differing by exactly $1$, so that the two roots define the same point of
$X_0(7)^+$ although they are different vectors with different (isometric)
complements.
\end{exercise}

\section*{Chapter note}

Theorems of Morrison~\cite{Morrison1984} and Shioda--Inose~\cite{ShiodaInose1977}
are quoted; the reader should check the normalization of the Hodge isometry in
Morrison's definition against the source. The tables of walls are elementary
computations reproduced for this book; the lattice halves are kernel-checked as
stated. The observation about ampleness follows Dolgachev's Thm.~7.3 and the
inequality on discriminants is immediate from Proposition~\ref{prop:index}.
